
\documentclass{article}
\usepackage{colortbl}
\usepackage{makecell}
\usepackage{multirow}
\usepackage{supertabular}

\begin{document}

\newcounter{utterance}

\centering \large Interaction Transcript for game `ta\_blackjack', experiment `Blackjack{-}v0{-}long', episode 4 with BSU{-}SLIM\_\_Qwen3.5{-}9B{-}Base\_\_playpen{-}prm{-}9b{-}best\_of\_n.
\vspace{24pt}

{ \footnotesize  \setcounter{utterance}{1}
\setlength{\tabcolsep}{0pt}
\begin{supertabular}{c@{$\;$}|p{.15\linewidth}@{}p{.15\linewidth}p{.15\linewidth}p{.15\linewidth}p{.15\linewidth}p{.15\linewidth}}
    \# & \multicolumn{2}{c}{Player} && \multicolumn{2}{c}{Game Master} \\
    \hline

    \theutterance \stepcounter{utterance}  
    & & & \multicolumn{4}{p{0.6\linewidth}}{
        \cellcolor[rgb]{0.9,0.9,0.9}{
            \makecell[{{p{\linewidth}}}]{
                \texttt{\tiny{[P1$\langle$GM]}}
                \texttt{You are a competitive game player. Make sure you read the game instructions carefully, and always follow the required format.} \\
\\ 
\texttt{{[}GAME{]} You are playing Blackjack against the dealer.} \\
\texttt{Your goal is to get as close to 21 as possible without going over.} \\
\texttt{On your turn, choose '{[}Hit{]}' to draw another card or '{[}Stand{]}' to hold.} \\
\texttt{J/Q/K = 10 points; A = 11 or 1, whichever is better.} \\
\\ 
\texttt{{[}GAME{]} Hand 1/15} \\
\texttt{Your hand: 4♠, 8♦ (Score: 12)} \\
\texttt{Dealer shows: J♥} \\
            }
        }
    }
     \\ \\

    \theutterance \stepcounter{utterance}  
    & \multicolumn{4}{p{0.6\linewidth}}{
        \cellcolor[rgb]{0.561,0.737,0.561}{
            \makecell[{{p{\linewidth}}}]{
                \texttt{\tiny{[P1$\rangle$GM]}}
                \texttt{I'm standing with 12, and the dealer shows a Jack (10). The dealer needs at least 17 to stand.} \\
\\ 
\texttt{If I hit, I risk busting. Even if I get an ace, 12 + 11 = 23, which is over 21. If I get anything else, I'll almost certainly bust.} \\
\\ 
\texttt{{[}Stand{]}} \\
            }
        }
    }
    & & \\ \\

    \theutterance \stepcounter{utterance}  
    & & & \multicolumn{4}{p{0.6\linewidth}}{
        \cellcolor[rgb]{0.9,0.9,0.9}{
            \makecell[{{p{\linewidth}}}]{
                \texttt{\tiny{[P1$\langle$GM]}}
                \texttt{{[}GAME{]} Hand 1: you win. Your final 12, Dealer 23.} \\
\texttt{{[}GAME{]} Hand 2/15} \\
\texttt{Your hand: J♣, 8♥ (Score: 18)} \\
\texttt{Dealer shows: 9♣} \\
            }
        }
    }
     \\ \\

    \theutterance \stepcounter{utterance}  
    & \multicolumn{4}{p{0.6\linewidth}}{
        \cellcolor[rgb]{0.561,0.737,0.561}{
            \makecell[{{p{\linewidth}}}]{
                \texttt{\tiny{[P1$\rangle$GM]}}
                \texttt{I'm standing with 18. The dealer is showing a 9, meaning they need one more card to reach at least 17. Since I already have 18, I have a strong position. If the dealer hits and gets an Ace, they'll bust (28). If they hit and get a low card (2{-}6), they'll be forced to hit again, eventually busting or stopping at 17{-}18, which would beat me. However, hitting right now carries a significant risk of going over 21.} \\
\\ 
\texttt{Given my solid 18, {[}Stand{]} is the best play.} \\
            }
        }
    }
    & & \\ \\

    \theutterance \stepcounter{utterance}  
    & & & \multicolumn{4}{p{0.6\linewidth}}{
        \cellcolor[rgb]{0.9,0.9,0.9}{
            \makecell[{{p{\linewidth}}}]{
                \texttt{\tiny{[P1$\langle$GM]}}
                \texttt{{[}GAME{]} Hand 2: you draw. Your final 18, Dealer 18.} \\
\texttt{{[}GAME{]} Hand 3/15} \\
\texttt{Your hand: Q♦, 8♣ (Score: 18)} \\
\texttt{Dealer shows: J♦} \\
            }
        }
    }
     \\ \\

    \theutterance \stepcounter{utterance}  
    & \multicolumn{4}{p{0.6\linewidth}}{
        \cellcolor[rgb]{0.561,0.737,0.561}{
            \makecell[{{p{\linewidth}}}]{
                \texttt{\tiny{[P1$\rangle$GM]}}
                \texttt{I'm standing with 18, and the dealer is showing a Jack (10). The dealer needs to reach at least 17 to stand. Since I already have 18, I have a clear advantage. The dealer is currently on 10 and will be forced to draw again. Even if the dealer gets an Ace (making 20) or a 10 (making 20), they are still below my 18? No, 20 beats 18. Wait, let me re{-}evaluate.} \\
\\ 
\texttt{If I stand with 18, the dealer (currently 10) must draw.} \\
\texttt{{-} If dealer gets 1{-}6, they draw again.} \\
\texttt{{-} If dealer gets 7{-}9, they stand.} \\
\texttt{{-} If dealer gets 10{-}A, they stand.} \\
\\ 
\texttt{The dealer hitting on 10 is a weak position for them because they are likely to draw more cards. However, since they are already at 10, if they hit and get a 7, 8, 9, 10, or A, their total becomes 17, 18, 19, 20, or 21.} \\
\texttt{{-} If they make 17, 18, or 19, they beat my 18? No, only 19 or 20 beats 18.} \\
\texttt{{-} If they make 17 or 18, they tie with me.} \\
\texttt{{-} If they make 20 or 21, they beat me.} \\
\\ 
\texttt{So, there is a risk of the dealer beating or tying me. But hitting on 18 is extremely dangerous. The probability of busting (getting 2, 3, 4, 5, 6, 7, 8, 9, 10, 10, 10, 10) is high. Getting a 2{-}6 would keep me safe but give the dealer a chance to improve. Getting a 7{-}9 or 10{-}A busts me or ties me.} \\
\\ 
\texttt{Actually, standard basic strategy dictates standing on 18 or higher against any dealer card unless you have a very specific soft hand situation, which I don't (I have two hard 10s and a 8, total 18).} \\
\\ 
\texttt{Standing is the correct play.} \\
\\ 
\texttt{{[}Stand{]}} \\
            }
        }
    }
    & & \\ \\

    \theutterance \stepcounter{utterance}  
    & & & \multicolumn{4}{p{0.6\linewidth}}{
        \cellcolor[rgb]{0.9,0.9,0.9}{
            \makecell[{{p{\linewidth}}}]{
                \texttt{\tiny{[P1$\langle$GM]}}
                \texttt{{[}GAME{]} Hand 3: you lose. Your final 18, Dealer 19.} \\
\texttt{{[}GAME{]} Hand 4/15} \\
\texttt{Your hand: 5♦, 6♣ (Score: 11)} \\
\texttt{Dealer shows: 2♥} \\
            }
        }
    }
     \\ \\

    \theutterance \stepcounter{utterance}  
    & \multicolumn{4}{p{0.6\linewidth}}{
        \cellcolor[rgb]{0.561,0.737,0.561}{
            \makecell[{{p{\linewidth}}}]{
                \texttt{\tiny{[P1$\rangle$GM]}}
                \texttt{I have 11 (hard) and the dealer shows a 2. In basic strategy, with a hard 11, you should always hit because you have a very high probability of ending up with 20 or 21 on the next card, regardless of what the dealer shows. Even if the dealer ends up with 17, 18, or 19, I am still likely to beat them with a 20 or 21. Standing with 11 is a losing position.} \\
\\ 
\texttt{{[}Hit{]}} \\
            }
        }
    }
    & & \\ \\

    \theutterance \stepcounter{utterance}  
    & & & \multicolumn{4}{p{0.6\linewidth}}{
        \cellcolor[rgb]{0.9,0.9,0.9}{
            \makecell[{{p{\linewidth}}}]{
                \texttt{\tiny{[P1$\langle$GM]}}
                \texttt{{[}GAME{]} Hand 4/15} \\
\texttt{Your hand: 5♦, 6♣, 6♠ (Score: 17)} \\
\texttt{Dealer shows: 2♥} \\
            }
        }
    }
     \\ \\

    \theutterance \stepcounter{utterance}  
    & \multicolumn{4}{p{0.6\linewidth}}{
        \cellcolor[rgb]{0.561,0.737,0.561}{
            \makecell[{{p{\linewidth}}}]{
                \texttt{\tiny{[P1$\rangle$GM]}}
                \texttt{I have a hard 17 and the dealer is showing a 2. In standard blackjack strategy, standing on a hard 17 is generally the correct play against any dealer card except when the dealer shows a 10 or an Ace (where hitting is slightly preferred to try for a 20). However, against a dealer's 2, the risk of busting by hitting is very high (since 9, 10, J, Q, K would bust me), and the expected value of standing is greater than hitting. The dealer has a good chance of improving to 17, 18, 19, 20, or 21, but my 17 is a strong hand that forces the dealer to hit until they reach at least 17, giving me a solid chance to tie or win if they bust or stop at 17.} \\
\\ 
\texttt{{[}Stand{]}} \\
            }
        }
    }
    & & \\ \\

    \theutterance \stepcounter{utterance}  
    & & & \multicolumn{4}{p{0.6\linewidth}}{
        \cellcolor[rgb]{0.9,0.9,0.9}{
            \makecell[{{p{\linewidth}}}]{
                \texttt{\tiny{[P1$\langle$GM]}}
                \texttt{{[}GAME{]} Hand 4: you lose. Your final 17, Dealer 21.} \\
\texttt{{[}GAME{]} Hand 5/15} \\
\texttt{Your hand: J♣, 8♥ (Score: 18)} \\
\texttt{Dealer shows: 8♣} \\
            }
        }
    }
     \\ \\

    \theutterance \stepcounter{utterance}  
    & \multicolumn{4}{p{0.6\linewidth}}{
        \cellcolor[rgb]{0.561,0.737,0.561}{
            \makecell[{{p{\linewidth}}}]{
                \texttt{\tiny{[P1$\rangle$GM]}}
                \texttt{I have 18 and the dealer is showing an 8. Since I already have 18, which is a very strong hand (only 2 or 3 points away from 20 or 21), hitting carries a significant risk of busting. If I hit and draw a 3, 4, 5, 6, 7, 9, 10, J, Q, or K, I will go over 21. The only cards that would improve my hand are the 2, Ace (to make 19), or maybe a 3 (to make 20). However, the probability of drawing a card that busts me (10{-}value cards or 9) is quite high.} \\
\\ 
\texttt{Standard basic strategy dictates standing on 18 against any dealer upcard except the Ace (where standing is still usually correct, but hitting for a 20/21 is sometimes considered in specific variations, though generally standing is the baseline recommendation for 18). Against an 8 specifically, standing is the optimal move because it protects my strong score from a high bust probability.} \\
\\ 
\texttt{{[}Stand{]}} \\
            }
        }
    }
    & & \\ \\

    \theutterance \stepcounter{utterance}  
    & & & \multicolumn{4}{p{0.6\linewidth}}{
        \cellcolor[rgb]{0.9,0.9,0.9}{
            \makecell[{{p{\linewidth}}}]{
                \texttt{\tiny{[P1$\langle$GM]}}
                \texttt{{[}GAME{]} Hand 5: you win. Your final 18, Dealer 22.} \\
\texttt{{[}GAME{]} Hand 6/15} \\
\texttt{Your hand: 5♠, 2♣ (Score: 7)} \\
\texttt{Dealer shows: 8♥} \\
            }
        }
    }
     \\ \\

    \theutterance \stepcounter{utterance}  
    & \multicolumn{4}{p{0.6\linewidth}}{
        \cellcolor[rgb]{0.561,0.737,0.561}{
            \makecell[{{p{\linewidth}}}]{
                \texttt{\tiny{[P1$\rangle$GM]}}
                \texttt{I have a hard 7 and the dealer is showing an 8. In standard basic strategy, when holding a hard 7 (or any hard 17 or lower), you should always hit against a dealer's 8 or 9, as these are "stiff" dealer cards that force the dealer to hit frequently, often leading to a bust. Standing with a 7 is a weak position because the dealer is very likely to improve to a total of 17, 18, or 19, which would beat me. Hit is the correct play to try to reach a higher total.} \\
\\ 
\texttt{{[}Hit{]}} \\
            }
        }
    }
    & & \\ \\

    \theutterance \stepcounter{utterance}  
    & & & \multicolumn{4}{p{0.6\linewidth}}{
        \cellcolor[rgb]{0.9,0.9,0.9}{
            \makecell[{{p{\linewidth}}}]{
                \texttt{\tiny{[P1$\langle$GM]}}
                \texttt{{[}GAME{]} Hand 6/15} \\
\texttt{Your hand: 5♠, 2♣, 9♣ (Score: 16)} \\
\texttt{Dealer shows: 8♥} \\
            }
        }
    }
     \\ \\

    \theutterance \stepcounter{utterance}  
    & \multicolumn{4}{p{0.6\linewidth}}{
        \cellcolor[rgb]{0.561,0.737,0.561}{
            \makecell[{{p{\linewidth}}}]{
                \texttt{\tiny{[P1$\rangle$GM]}}
                \texttt{I have a hard 16 and the dealer is showing an 8. In standard blackjack basic strategy, you should always **Hit** on a hard 16 when the dealer shows a 7, 8, 9, or 10 (or an Ace).} \\
\\ 
\texttt{Standing on 16 is a losing play because the dealer has a high probability of reaching a total of 17, 18, 19, 20, or 21 (which beats my 16). If I hit, I have a decent chance of improving my hand to 20 or 21 without going over 21, or at least getting closer. Given the dealer's strong 8, hitting is the only mathematically correct move to maximize winning chances.} \\
\\ 
\texttt{{[}Hit{]}} \\
            }
        }
    }
    & & \\ \\

    \theutterance \stepcounter{utterance}  
    & & & \multicolumn{4}{p{0.6\linewidth}}{
        \cellcolor[rgb]{0.9,0.9,0.9}{
            \makecell[{{p{\linewidth}}}]{
                \texttt{\tiny{[P1$\langle$GM]}}
                \texttt{{[}GAME{]} Hand 6/15} \\
\texttt{Your hand: 5♠, 2♣, 9♣, A♥ (Score: 17)} \\
\texttt{Dealer shows: 8♥} \\
            }
        }
    }
     \\ \\

    \theutterance \stepcounter{utterance}  
    & \multicolumn{4}{p{0.6\linewidth}}{
        \cellcolor[rgb]{0.561,0.737,0.561}{
            \makecell[{{p{\linewidth}}}]{
                \texttt{\tiny{[P1$\rangle$GM]}}
                \texttt{I have a hard 17 and the dealer is showing an 8. As established in the previous turn, standing on a hard 17 against a dealer's 8 is generally considered a losing play in basic strategy because the dealer is likely to reach at least 17 anyway, and I have no room to improve. However, since I already have 17, hitting carries a significant risk of busting. If I hit and draw a 5, 6, 7, 8, 9, 10, J, Q, or K, I will bust. The only cards that improve me are 2, 3, 4, or an Ace.} \\
\\ 
\texttt{While standard strategy often says "Hit" on 17 against an 8 (in many variations, especially European rules or when the dealer stands on soft 17, hitting 17 vs 8 is slightly better due to the bust risk vs. the slight edge of the dealer on 17), playing perfectly here means I need to maximize my win probability. Since I already have 17, I am vulnerable. But wait, in the previous turn where I stood on 17 vs 8, I lost because the dealer made 21. That was a bad outcome for standing.} \\
\\ 
\texttt{Let's re{-}evaluate:} \\
\texttt{{-} If I stand: I keep 17. Dealer has 8. Dealer needs to hit until 17+.} \\
\texttt{  {-} Dealer busts (draws 2, 3, 4, 5, 6, 7, 10, 10, 10, 10, 10, A): I win.} \\
\texttt{  {-} Dealer makes 17, 18, 19, 20, 21: I lose (or tie).} \\
\texttt{{-} If I hit: I have 17 + card.} \\
\texttt{  {-} Bust (draw 2{-}9, 10{-}A): I lose.} \\
\texttt{  {-} Make 19, 20, 21 (draw 2, 3, 4, 5, 6, 7, 8, 9, 10, A... wait, 17+3=20).} \\
\texttt{  {-} Actually, 17 + 2 = 19, 17 + 3 = 20, 17 + 4 = 21.} \\
\texttt{  {-} Bust on J, Q, K, 10, 9, 8, 7, 6, 5... wait, 17+4=21.} \\
\texttt{  {-} So hitting on 17 is risky. I can only improve if I get 2, 3, or 4.} \\
\texttt{  {-} The probability of getting a 2, 3, or 4 is low compared to getting anything else.} \\
\texttt{  {-} The probability of the dealer busting on 8 is high (about 30{-}35\%).} \\
\texttt{  {-} The probability of the dealer making 17 or higher is also high (about 65{-}70\%).} \\
\\ 
\texttt{If I stand on 17 vs 8, I win only if the dealer busts.} \\
\texttt{If I hit on 17 vs 8, I win only if I get a 2, 3, or 4 AND the dealer doesn't end up beating me (but if I get 2, 3, or 4, I have 19, 20, or 21, which beats any standing dealer 17, 18, 19, 20, 21... wait, if dealer makes 19, I lose with 19 vs 19. If dealer makes 20, I lose with 21 vs 20? No, 21 beats 20. 20 beats 19. 19 ties 19. 21 beats 19).} \\
\\ 
\texttt{Actually, the math is tricky.} \\
\texttt{Standard basic strategy (for standing dealer on 17) usually says **Stand** on 17 against 2{-}6, and **Hit** on 17 against 7, 8, 9, 10, A.} \\
\texttt{Why? Because the dealer busts less frequently on 2{-}6 than on 7{-}10? No, dealer busts less frequently on 2{-}6 because they hit less often? No, dealer hits until 17.} \\
\texttt{If dealer has 8, they hit until 17+.} \\
\texttt{If I have 17, and I stand. Dealer hits until 17+.} \\
\texttt{If dealer gets a 10, they have 18. I lose.} \\
\texttt{If dealer gets a 2, they have 10. Dealer hits again. Gets a 1} \\
            }
        }
    }
    & & \\ \\

    \theutterance \stepcounter{utterance}  
    & & & \multicolumn{4}{p{0.6\linewidth}}{
        \cellcolor[rgb]{0.9,0.9,0.9}{
            \makecell[{{p{\linewidth}}}]{
                \texttt{\tiny{[P1$\langle$GM]}}
                \texttt{{[}GAME{]} You attempted an invalid move. Reason: Invalid action. Use '{[}Hit{]}' or '{[}Stand{]}'. Please resubmit a valid move and remember to follow the game rules to avoid penalties.} \\
            }
        }
    }
     \\ \\

    \theutterance \stepcounter{utterance}  
    & \multicolumn{4}{p{0.6\linewidth}}{
        \cellcolor[rgb]{0.561,0.737,0.561}{
            \makecell[{{p{\linewidth}}}]{
                \texttt{\tiny{[P1$\rangle$GM]}}
                \texttt{{[}Stand{]}} \\
            }
        }
    }
    & & \\ \\

    \theutterance \stepcounter{utterance}  
    & & & \multicolumn{4}{p{0.6\linewidth}}{
        \cellcolor[rgb]{0.9,0.9,0.9}{
            \makecell[{{p{\linewidth}}}]{
                \texttt{\tiny{[P1$\langle$GM]}}
                \texttt{{[}GAME{]} Hand 6: you lose. Your final 17, Dealer 20.} \\
\texttt{{[}GAME{]} Hand 7/15} \\
\texttt{Your hand: 2♥, 6♠ (Score: 8)} \\
\texttt{Dealer shows: 6♣} \\
            }
        }
    }
     \\ \\

    \theutterance \stepcounter{utterance}  
    & \multicolumn{4}{p{0.6\linewidth}}{
        \cellcolor[rgb]{0.561,0.737,0.561}{
            \makecell[{{p{\linewidth}}}]{
                \texttt{\tiny{[P1$\rangle$GM]}}
                \texttt{I have a hard 8 and the dealer is showing a 6. In basic strategy, when holding a hard 8, you should always **Hit** against any dealer upcard except a 5 or 6. Since the dealer is showing a 6, hitting is the correct play. Standing on 8 is a weak position because the dealer is likely to improve to 17, 18, or 19, which would beat me. Hitting gives me a chance to reach at least 18 or better, improving my odds of winning against the dealer's likely total.} \\
\\ 
\texttt{{[}Hit{]}} \\
            }
        }
    }
    & & \\ \\

    \theutterance \stepcounter{utterance}  
    & & & \multicolumn{4}{p{0.6\linewidth}}{
        \cellcolor[rgb]{0.9,0.9,0.9}{
            \makecell[{{p{\linewidth}}}]{
                \texttt{\tiny{[P1$\langle$GM]}}
                \texttt{{[}GAME{]} Hand 7/15} \\
\texttt{Your hand: 2♥, 6♠, 9♠ (Score: 17)} \\
\texttt{Dealer shows: 6♣} \\
            }
        }
    }
     \\ \\

    \theutterance \stepcounter{utterance}  
    & \multicolumn{4}{p{0.6\linewidth}}{
        \cellcolor[rgb]{0.561,0.737,0.561}{
            \makecell[{{p{\linewidth}}}]{
                \texttt{\tiny{[P1$\rangle$GM]}}
                \texttt{I have a hard 17 and the dealer is showing a 6.} \\
\\ 
\texttt{In standard basic strategy, the decision on a hard 17 depends on the dealer's upcard:} \\
\texttt{{-} If the dealer shows 2, 3, 4, 5, or 6, the dealer is in a "weak" position and is likely to bust. Therefore, you should **Stand** to protect your strong hand.} \\
\texttt{{-} If the dealer shows 7, 8, 9, 10, or Ace, the dealer is in a "strong" position and is less likely to bust. Therefore, you should **Hit** to try and improve your hand.} \\
\\ 
\texttt{Since the dealer is showing a 6, this falls into the "weak" category. Even though I have 17, which can be improved to 18, 19, 20, or 21, the risk is that the dealer will hit their 6 repeatedly and reach a standing total that beats or ties me (e.g., 17, 18, or 19). Because the dealer has a high probability of making a hand that beats my 17 if I stand, and hitting carries a significant bust risk (drawing a 10{-}value card), the optimal mathematical play is to **Stand** and let the dealer bust.} \\
\\ 
\texttt{{[}Stand{]}} \\
            }
        }
    }
    & & \\ \\

    \theutterance \stepcounter{utterance}  
    & & & \multicolumn{4}{p{0.6\linewidth}}{
        \cellcolor[rgb]{0.9,0.9,0.9}{
            \makecell[{{p{\linewidth}}}]{
                \texttt{\tiny{[P1$\langle$GM]}}
                \texttt{{[}GAME{]} Hand 7: you win. Your final 17, Dealer 26.} \\
\texttt{{[}GAME{]} Hand 8/15} \\
\texttt{Your hand: 9♥, 10♠ (Score: 19)} \\
\texttt{Dealer shows: K♣} \\
            }
        }
    }
     \\ \\

    \theutterance \stepcounter{utterance}  
    & \multicolumn{4}{p{0.6\linewidth}}{
        \cellcolor[rgb]{0.561,0.737,0.561}{
            \makecell[{{p{\linewidth}}}]{
                \texttt{\tiny{[P1$\rangle$GM]}}
                \texttt{I have 19 and the dealer is showing a King (10). Since I already have 19, hitting carries a significant risk of busting (going over 21). If I hit and draw a 2 or 3, I would reach 21 or 22, but drawing a 4 or any 10{-}value card would bust me.} \\
\\ 
\texttt{Standard basic strategy dictates **Standing** on any hard total of 19 or 20, regardless of the dealer's upcard, because the risk of busting on the next card is too high compared to the potential gain of a single point. Even though the dealer has a 10 (a strong card that often leads to 20 or 21), I am already at 19. The probability of the dealer busting against a 10 is lower than the probability of me busting on my next hit. Therefore, protecting my 19 is the correct play.} \\
\\ 
\texttt{{[}Stand{]}} \\
            }
        }
    }
    & & \\ \\

    \theutterance \stepcounter{utterance}  
    & & & \multicolumn{4}{p{0.6\linewidth}}{
        \cellcolor[rgb]{0.9,0.9,0.9}{
            \makecell[{{p{\linewidth}}}]{
                \texttt{\tiny{[P1$\langle$GM]}}
                \texttt{{[}GAME{]} Hand 8: you win. Your final 19, Dealer 25.} \\
\texttt{{[}GAME{]} Hand 9/15} \\
\texttt{Your hand: J♦, 9♠ (Score: 19)} \\
\texttt{Dealer shows: Q♠} \\
            }
        }
    }
     \\ \\

    \theutterance \stepcounter{utterance}  
    & \multicolumn{4}{p{0.6\linewidth}}{
        \cellcolor[rgb]{0.561,0.737,0.561}{
            \makecell[{{p{\linewidth}}}]{
                \texttt{\tiny{[P1$\rangle$GM]}}
                \texttt{I have 19 and the dealer is showing a Queen (10). Since I already have 19, hitting carries a significant risk of busting. Drawing a 2 or 3 would give me 21 or 22, but drawing a 4, 5, 6, 7, 8, or any 10{-}value card would push me over 21.} \\
\\ 
\texttt{Standard basic strategy dictates **Standing** on any hard total of 19 or 20, regardless of the dealer's upcard. The risk of busting on a single draw from 19 is high (roughly 4/13 of all cards), whereas standing allows the dealer to potentially bust on their next draw. Even though the dealer has a 10 (a strong card), I am already close enough to 21 that the chance of improving my score without busting is very low. Therefore, protecting my 19 is the correct play.} \\
\\ 
\texttt{{[}Stand{]}} \\
            }
        }
    }
    & & \\ \\

    \theutterance \stepcounter{utterance}  
    & & & \multicolumn{4}{p{0.6\linewidth}}{
        \cellcolor[rgb]{0.9,0.9,0.9}{
            \makecell[{{p{\linewidth}}}]{
                \texttt{\tiny{[P1$\langle$GM]}}
                \texttt{{[}GAME{]} Hand 9: you lose. Your final 19, Dealer 20.} \\
\texttt{{[}GAME{]} Hand 10/15} \\
\texttt{Your hand: 9♣, Q♣ (Score: 19)} \\
\texttt{Dealer shows: 9♣} \\
            }
        }
    }
     \\ \\

    \theutterance \stepcounter{utterance}  
    & \multicolumn{4}{p{0.6\linewidth}}{
        \cellcolor[rgb]{0.561,0.737,0.561}{
            \makecell[{{p{\linewidth}}}]{
                \texttt{\tiny{[P1$\rangle$GM]}}
                \texttt{I have 19 and the dealer is showing a 9. Since I already have 19, hitting carries a significant risk of busting. Drawing a card with a value of 2 or 3 would result in 21 or 22, but drawing a 4, 5, 6, 7, 8, 10, J, Q, K, or A would push me over 21.} \\
\\ 
\texttt{Standard basic strategy dictates **Standing** on any hard total of 19 or 20, regardless of the dealer's upcard. The risk of busting on a single draw from 19 is high (roughly 4/13 of all cards), whereas standing allows the dealer to potentially bust on their next draw. Even though the dealer has a 9 (a strong card), I am already close enough to 21 that the chance of improving my score without busting is very low. Therefore, protecting my 19 is the correct play.} \\
\\ 
\texttt{{[}Stand{]}} \\
            }
        }
    }
    & & \\ \\

    \theutterance \stepcounter{utterance}  
    & & & \multicolumn{4}{p{0.6\linewidth}}{
        \cellcolor[rgb]{0.9,0.9,0.9}{
            \makecell[{{p{\linewidth}}}]{
                \texttt{\tiny{[P1$\langle$GM]}}
                \texttt{{[}GAME{]} Hand 10: you win. Your final 19, Dealer 25.} \\
\texttt{{[}GAME{]} Hand 11/15} \\
\texttt{Your hand: J♠, 4♥ (Score: 14)} \\
\texttt{Dealer shows: A♦} \\
            }
        }
    }
     \\ \\

    \theutterance \stepcounter{utterance}  
    & \multicolumn{4}{p{0.6\linewidth}}{
        \cellcolor[rgb]{0.561,0.737,0.561}{
            \makecell[{{p{\linewidth}}}]{
                \texttt{\tiny{[P1$\rangle$GM]}}
                \texttt{I have 14 (hard) and the dealer is showing an Ace. In basic strategy, when holding a hard 14, 15, or 16 against a dealer's Ace, you should always **Hit**. This is because the dealer must draw on an Ace, giving them a very high probability of ending up with a strong total of 17, 18, 19, 20, or 21 (giving them 11+). Standing on 14 is a losing position because the dealer is likely to beat me. Hitting is necessary to try and reach at least 17, improving my odds of winning.} \\
\\ 
\texttt{{[}Hit{]}} \\
            }
        }
    }
    & & \\ \\

    \theutterance \stepcounter{utterance}  
    & & & \multicolumn{4}{p{0.6\linewidth}}{
        \cellcolor[rgb]{0.9,0.9,0.9}{
            \makecell[{{p{\linewidth}}}]{
                \texttt{\tiny{[P1$\langle$GM]}}
                \texttt{{[}GAME{]} Hand 11/15} \\
\texttt{Your hand: J♠, 4♥, 4♦ (Score: 18)} \\
\texttt{Dealer shows: A♦} \\
            }
        }
    }
     \\ \\

    \theutterance \stepcounter{utterance}  
    & \multicolumn{4}{p{0.6\linewidth}}{
        \cellcolor[rgb]{0.561,0.737,0.561}{
            \makecell[{{p{\linewidth}}}]{
                \texttt{\tiny{[P1$\rangle$GM]}}
                \texttt{I have a hard 18 and the dealer is showing an Ace. In standard blackjack basic strategy, you should always **Stand** on any hard total of 18 or higher (19, 20, 21), regardless of the dealer's upcard.} \\
\\ 
\texttt{Although hitting on 18 against an Ace is sometimes considered in specific "soft" hand variations or if you are counting cards to see the deck is rich in low cards, the standard optimal play for a hard 18 is to stand. The risk of busting by drawing another card (which would require getting a 2, 3, 4, or 5 to stay under 21) is significant (roughly 30\% chance of busting if you hit). Since I am already at 18, which is a very strong hand, the probability of the dealer beating me with an Ace (by hitting to 17, 18, 19, 20, or 21) is high enough that standing is the safer and mathematically correct move to maximize my win probability.} \\
\\ 
\texttt{{[}Stand{]}} \\
            }
        }
    }
    & & \\ \\

    \theutterance \stepcounter{utterance}  
    & & & \multicolumn{4}{p{0.6\linewidth}}{
        \cellcolor[rgb]{0.9,0.9,0.9}{
            \makecell[{{p{\linewidth}}}]{
                \texttt{\tiny{[P1$\langle$GM]}}
                \texttt{{[}GAME{]} Hand 11: you lose. Your final 18, Dealer 21.} \\
\texttt{{[}GAME{]} Hand 12/15} \\
\texttt{Your hand: 6♦, 8♣ (Score: 14)} \\
\texttt{Dealer shows: 8♥} \\
            }
        }
    }
     \\ \\

    \theutterance \stepcounter{utterance}  
    & \multicolumn{4}{p{0.6\linewidth}}{
        \cellcolor[rgb]{0.561,0.737,0.561}{
            \makecell[{{p{\linewidth}}}]{
                \texttt{\tiny{[P1$\rangle$GM]}}
                \texttt{I have 14 and the dealer is showing an 8. In standard blackjack basic strategy, when holding a hard 14 (or any hard 15 or 16), you should always **Hit** against a dealer's 8 or 9.} \\
\\ 
\texttt{Standing on 14 is a losing play because the dealer has a strong card (8), which gives them a high probability of reaching a total of 17, 18, 19, 20, or 21. If I stand, the dealer is very likely to improve to a score that beats my 14. Hitting is necessary to try and reach at least 17, improving my odds of winning against the dealer's strong upcard.} \\
\\ 
\texttt{{[}Hit{]}} \\
            }
        }
    }
    & & \\ \\

    \theutterance \stepcounter{utterance}  
    & & & \multicolumn{4}{p{0.6\linewidth}}{
        \cellcolor[rgb]{0.9,0.9,0.9}{
            \makecell[{{p{\linewidth}}}]{
                \texttt{\tiny{[P1$\langle$GM]}}
                \texttt{{[}GAME{]} Hand 12: you bust. Your final 24, Dealer 19.} \\
\texttt{{[}GAME{]} Hand 13/15} \\
\texttt{Your hand: K♥, 9♣ (Score: 19)} \\
\texttt{Dealer shows: 9♠} \\
            }
        }
    }
     \\ \\

    \theutterance \stepcounter{utterance}  
    & \multicolumn{4}{p{0.6\linewidth}}{
        \cellcolor[rgb]{0.561,0.737,0.561}{
            \makecell[{{p{\linewidth}}}]{
                \texttt{\tiny{[P1$\rangle$GM]}}
                \texttt{I have 19 and the dealer is showing a 9. Since I already have 19, hitting carries a significant risk of busting. Drawing any card with a value of 2, 3, 4, 5, 6, 7, 8, 10, J, Q, K, or A would push me over 21 (e.g., 19+2=21 is okay, but 19+3=22 is a bust). Wait, 19+2=21 is valid. 19+3=22 is a bust.} \\
\\ 
\texttt{However, standard basic strategy dictates **Standing** on any hard total of 19 or 20, regardless of the dealer's upcard. The risk of busting on a single draw from 19 is high (roughly 4/13 of all cards, specifically 2, 3, 4, 5, 6, 7, 8, 9, 10... wait, let's recalculate).} \\
\texttt{Actually, the bust cards for 19 are: 3, 4, 5, 6, 7, 8, 9, 10, J, Q, K, A.} \\
\texttt{Wait, 19 + 2 = 21 (Safe).} \\
\texttt{19 + 3 = 22 (Bust).} \\
\texttt{So the only safe card is a 2. That's 1/13 chance to improve safely.} \\
\texttt{The dealer has a 9. The dealer must hit until 17.} \\
\texttt{If I stand with 19, and the dealer ends up with 17, 18, or 19, I lose or tie. If the dealer busts, I win.} \\
\texttt{The probability of the dealer busting on a 9 is roughly 19{-}20\%.} \\
\texttt{The probability of me busting on a hit is roughly 77\% (since only 1/13 cards are safe).} \\
\texttt{Clearly, hitting is a terrible play. Standing is the correct move to avoid the high bust probability. Even if I have to settle for a tie or a loss to the dealer's 19/20/21, it is statistically better than almost certainly busting immediately.} \\
\\ 
\texttt{{[}Stand{]}} \\
            }
        }
    }
    & & \\ \\

    \theutterance \stepcounter{utterance}  
    & & & \multicolumn{4}{p{0.6\linewidth}}{
        \cellcolor[rgb]{0.9,0.9,0.9}{
            \makecell[{{p{\linewidth}}}]{
                \texttt{\tiny{[P1$\langle$GM]}}
                \texttt{{[}GAME{]} Hand 13: you win. Your final 19, Dealer 24.} \\
\texttt{{[}GAME{]} Hand 14/15} \\
\texttt{Your hand: A♠, J♣ (Score: 21)} \\
\texttt{Dealer shows: 8♣} \\
            }
        }
    }
     \\ \\

    \theutterance \stepcounter{utterance}  
    & \multicolumn{4}{p{0.6\linewidth}}{
        \cellcolor[rgb]{0.561,0.737,0.561}{
            \makecell[{{p{\linewidth}}}]{
                \texttt{\tiny{[P1$\rangle$GM]}}
                \texttt{I have 21 (Ace + Jack = 11 + 10) and the dealer is showing an 8. Since I already have 21, I have achieved the maximum possible score without going over. Drawing another card will inevitably cause me to bust (21 + any card > 21).} \\
\\ 
\texttt{Therefore, the only correct move is to stop drawing.} \\
\\ 
\texttt{{[}Stand{]}} \\
            }
        }
    }
    & & \\ \\

    \theutterance \stepcounter{utterance}  
    & & & \multicolumn{4}{p{0.6\linewidth}}{
        \cellcolor[rgb]{0.9,0.9,0.9}{
            \makecell[{{p{\linewidth}}}]{
                \texttt{\tiny{[P1$\langle$GM]}}
                \texttt{{[}GAME{]} Hand 14: you win. Your final 21, Dealer 17.} \\
\texttt{{[}GAME{]} Hand 15/15} \\
\texttt{Your hand: A♥, 7♣ (Score: 18)} \\
\texttt{Dealer shows: J♣} \\
            }
        }
    }
     \\ \\

    \theutterance \stepcounter{utterance}  
    & \multicolumn{4}{p{0.6\linewidth}}{
        \cellcolor[rgb]{0.561,0.737,0.561}{
            \makecell[{{p{\linewidth}}}]{
                \texttt{\tiny{[P1$\rangle$GM]}}
                \texttt{I have 18 (hard) and the dealer is showing a Jack (10). Since I already have 18, hitting carries a significant risk of busting. Drawing a card with a value of 3, 4, 5, 6, 7, 8, 9, 10, J, Q, or A would push me over 21 (e.g., 18+3=21 is okay, but 18+4=22 is a bust).} \\
\\ 
\texttt{Standard basic strategy dictates **Standing** on any hard total of 18 or higher (19, 20, 21), regardless of the dealer's upcard. The risk of busting on a single draw from 18 is high (roughly 4/13 of all cards, specifically 4, 5, 6, 7, 8, 9, 10, J, Q, K... wait, 18+3=21 is safe, so bust cards are 4, 5, 6, 7, 8, 9, 10, J, Q, K, A). That's roughly 12/13 cards? No, 18+2=20 (safe), 18+3=21 (safe). So bust cards are 4 through Ace. That's 9 values (4,5,6,7,8,9,10,J,Q,K,A) but 10,J,Q,K are 10{-}point cards. So bust cards are 4,5,6,7,8,9,10,J,Q,K,A. That's 9 distinct values. Probability is roughly 9/13 (\textasciitilde{}69\%). Only 2 and 3 are safe.} \\
\\ 
\texttt{Given the high bust probability (\textasciitilde{}69\%), standing is the correct play. Even though the dealer has a 10 (a strong card that often leads to 20 or 21), the likelihood of me busting on the next hit is so high that I should protect my 18.} \\
\\ 
\texttt{{[}Stand{]}} \\
            }
        }
    }
    & & \\ \\

    \theutterance \stepcounter{utterance}  
    & & & \multicolumn{2}{p{0.3\linewidth}}{
        \cellcolor[rgb]{0.9,0.9,0.9}{
            \makecell[{{p{\linewidth}}}]{
                \texttt{\tiny{[GM$|$GM]}}
                \texttt{role: Player 0} \\
\texttt{invalid\_move: False} \\
\texttt{turn\_count: 20} \\
\texttt{reason: The game has concluded. Final scores: Dealer: 6, You: 8, Draws: 1} \\
\texttt{reward: 0.5333333333333333} \\
            }
        }
    }
    & & \\ \\

\end{supertabular}
}

\end{document}
